% =============================================================================
% Figure captions for: Membrane-Mediated Categorical Processing Architecture
% =============================================================================

\begin{figure*}[!htbp]
    \centering
    \includegraphics[width=\textwidth]{figures/panel_exp01_membrane_substrate.png}
    \caption{\textbf{Biological semiconductor characteristics of the amphiphilic bilayer substrate: diode behaviour, carrier transport, transistor switching, and tri-dimensional impedance.}
    (\textbf{A})~Current--voltage characteristic of the lipid bilayer P-N junction at physiological temperature $T = 310$~K, following the Shockley diode equation $I = I_0[\exp(eV / n\kB T) - 1]$ with saturation current $I_0 = 10^{-12}$~A and ideality factor $n = 1.5$. Forward bias shows the characteristic exponential onset; reverse bias exhibits sub-picoampere leakage. The built-in potential $V_{\mathrm{bi}} = 0.615$~V (dashed red) is derived from the carrier concentration ratio $\ln(p \cdot n / n_i^2)$ with $p = 2.80 \times 10^{12}$~cm$^{-3}$ (oscillatory holes) and $n = 1.12 \times 10^{12}$~cm$^{-3}$ (molecular oscillators), requiring zero fitted parameters.
    (\textbf{B})~Drift velocity $v_d$ of oscillatory hole carriers as a function of applied electric field $E$ on a double-logarithmic scale. At low fields, $v_d = \mu E$ with hole mobility $\mu = 0.0123$~cm$^2$\,V$^{-1}$\,s$^{-1}$ (blue dashed). At high fields ($E > 10^5$~V/m), the velocity saturates at $v_{\mathrm{sat}} \approx 80{,}000$~cm/s (red dashed) due to coupling-limited scattering in the lipid oscillator network. The crossover regime spans two decades of field strength, consistent with biological membrane operating conditions.
    (\textbf{C})~BMD transistor transfer characteristic: output current $I_{\mathrm{out}}$ as a function of the gate overlap integral $M = |\langle S_{\mathrm{in}} | P_{\mathrm{gate}} \rangle|^2 / (\langle S_{\mathrm{in}} | S_{\mathrm{in}} \rangle \langle P_{\mathrm{gate}} | P_{\mathrm{gate}} \rangle)$. Below the matching threshold $\theta = 0.5$, the transistor is in the off state ($I_{\mathrm{out}} \approx I_{\mathrm{in}} / G$). Above $\theta$, the transfer function $I_{\mathrm{out}} = G(M - \theta)/(1 - M + \theta) \cdot I_{\mathrm{in}}$ with gain $G = 1000$ produces the annotated on/off ratio of $42.1$, matching the theoretical prediction from the BMD filtering operator selectivity.
    (\textbf{D})~Three-dimensional surface of the R-C-L impedance magnitude $|Z(\omega, \Sk)| = |R(\Sk) + j\omega L(\Se) + 1/(j\omega C(\St))|$ as a function of angular frequency $\omega \in [10^3, 10^9]$~rad/s and knowledge entropy coordinate $\Sk \in [0, 1]$. The impedance is computed with $R(\Sk) = R_0 \exp(\Sk / \Sk^{\mathrm{ref}})$, $L = 3.14$~TH, $C = 318.3$~fF. The surface demonstrates simultaneous resistive, capacitive, and inductive character: the ridge at the L-C resonance frequency migrates with $\Sk$, confirming that the mode selection mechanism of the tri-dimensional R-C-L element is governed by the S-entropy coordinate, not by external bias.}
    \label{fig:membrane_substrate}
\end{figure*}

\begin{figure*}[!htbp]
    \centering
    \includegraphics[width=\textwidth]{figures/panel_exp02_categorical_apertures.png}
    \caption{\textbf{Categorical aperture filtering: trajectory selection, multipole field decay, ternary trisection, and progressive S-entropy refinement.}
    (\textbf{A})~Fifty particle trajectories (straight-line ballistic paths) incident on a circular categorical aperture of radius $R = 1$ in the $(x, y)$ configuration plane. Trajectories intersecting the aperture domain (blue) pass through unperturbed; those outside the domain (red) are blocked. Grey shading indicates the excluded region. The aperture exerts no force on transmitted particles ($W_{\mathrm{aperture}} = 0$ identically), performing geometric selection by position alone---confirming the Zero-Work Theorem (Theorem~\ref{thm:zero_work}).
    (\textbf{B})~Field strength decay for the three canonical aperture multipole orders on a double-logarithmic scale: monopole $|\mathbf{E}| \propto r^{-2}$ ($\ell = 0$, blue), dipole $|\mathbf{E}| \propto r^{-3}$ ($\ell = 1$, green), and quadrupole $|\mathbf{E}| \propto r^{-4}$ ($\ell = 2$, red). Distance $r$ spans $0.1$ to $10$ in reduced units. The steeper decay of higher-order multipoles reflects their more localised selectivity, with quadrupolar apertures (Hill coefficient $n_H = 4$) providing the sharpest spatial discrimination at the cost of shorter interaction range.
    (\textbf{C})~Ternary trisection of S-entropy space. Blue circles: $\log_3(V_k / V_0)$ showing the volume of the navigable region after $k$ successive ternary trisections ($k = 0, \ldots, 20$). Red diamonds: $\log_3(3^k)$, the number of distinguishable categorical states at depth $k$. Both scale linearly with $k$ at unit slope, confirming that each trisection simultaneously reduces the search volume by a factor of $3$ and increases the resolution by a factor of $3$. After $k = 20$ trisections, the addressable volume is $3^{-20} \approx 3 \times 10^{-10}$ of the original---sufficient for sub-nanometre precision in a $1$~cm system.
    (\textbf{D})~Three-dimensional scatter of $1000$ initial points uniformly distributed in the S-entropy cube $[0, 1]^3$, progressively filtered by five sequential categorical apertures. Points surviving all five stages (darkest colour) cluster in a small region of S-entropy space, demonstrating the multiplicative selectivity of cascaded apertures: total selectivity $\eta_{\mathrm{total}} = \prod_{i=1}^5 \eta_i$. Each stage removes approximately $2/3$ of the surviving population, consistent with the ternary trisection model.}
    \label{fig:categorical_apertures}
\end{figure*}

\begin{figure*}[!htbp]
    \centering
    \includegraphics[width=\textwidth]{figures/panel_exp03_trajectory_completion.png}
    \caption{\textbf{Trajectory completion dynamics: categorical speedup, backward navigation, and contraction mapping convergence.}
    (\textbf{A})~Operations required to navigate to a target state as a function of problem size $N$ on a double-logarithmic scale. Categorical navigation (red) requires $\lceil \log_3 N \rceil$ operations, growing logarithmically. Conventional sequential search at $O(N)$ (green, dashed) and comparison-based sorting at $O(N \log N)$ (blue, dashed) are overlaid for reference. At $N = 10^6$, the categorical approach requires $\sim 13$ operations versus $\sim 10^6$ sequential steps---a speedup of $\sim 79{,}000\times$. At $N = 10^{12}$, the speedup exceeds $10^{10}\times$.
    (\textbf{B})~Speedup factor $N / \log_3 N$ as a function of $N$ on a double-logarithmic scale. The speedup grows nearly linearly with $N$ (slope $\approx 1$ on log-log axes), confirming that the categorical advantage is not a constant factor but a complexity class separation: $O(\log N)$ versus $O(N)$.
    (\textbf{C})~Three-dimensional trajectory completion in S-entropy space $(\Sk, \St, \Se) \in [0, 1]^3$. Five trajectories (coloured lines) originate from distinct random starting points (circles) and converge toward a common target terminus (gold star) via successive ternary refinements. Intermediate waypoints (small dots) show the step-wise navigation. All trajectories reach the $\eps$-neighbourhood of the terminus regardless of initial position, confirming the Path-Independent Convergence Theorem.
    (\textbf{D})~Categorical distance $d_{\mathrm{cat}}(t)$ to the target terminus as a function of iteration step for the five trajectories in~(\textbf{C}), on a logarithmic $y$-axis. All five curves decay exponentially (linear on the semi-log plot), with contraction ratio $c \approx 1/3$ per step---consistent with the theoretical prediction from the ternary trisection model. Convergence to machine precision ($d_{\mathrm{cat}} < 10^{-15}$) is achieved within $15$--$18$ iterations for all starting points, confirming the Banach fixed-point convergence guarantee.}
    \label{fig:trajectory_completion}
\end{figure*}

\begin{figure*}[!htbp]
    \centering
    \includegraphics[width=\textwidth]{figures/panel_exp04_processing_cycle.png}
    \caption{\textbf{Processing cycle dynamics: decay envelope intersection, Kuramoto phase portrait, regime cycling, and R-C-L mode selection.}
    (\textbf{A})~Dual decay envelope model for the processing cycle. The input coupling envelope $P(t) = P_0 \exp(-t/\tau_1)$ with $\tau_1 = 50$~ms (blue) and the integration envelope $T(t) = T_0 (1 - e^{-t/\tau_2}) e^{-t/\tau_3}$ with $\tau_2 = 20$~ms (rise) and $\tau_3 = 100$~ms (decay) are plotted over $t \in [0, 500]$~ms. The shaded overlap region (purple) defines the processing window---the temporal interval during which both envelopes simultaneously exceed threshold. The window duration $\Delta t \approx 100$--$300$~ms is set entirely by the decay time constants without free parameters.
    (\textbf{B})~Kuramoto phase portrait: coupling strength $K$ versus phase variance $\sigma^2(\varphi)$. The equilibrium curve $\sigma^2_{\mathrm{eq}} = 1/\sqrt{K}$ (solid red) traces the thermodynamic variance floor. Simulated trajectories (grey curves) from five initial conditions spiral inward toward the equilibrium, confirming that the system is a gradient flow on the free energy $F = \kB T \cdot \sigma^2$. Green circles mark the equilibrium points at $K = 0.5, 1.0, 2.0, 3.0$.
    (\textbf{C})~Regime cycling: Kuramoto order parameter $R(t)$ over a $10$-second interval, oscillating through the five operational regimes. Coloured background bands demarcate: turbulent (red, $R < 0.3$), cascade (orange), coherent (green, $R > 0.8$), and phase-locked (blue, $R > 0.95$). The system executes a complete regime cycle---coherent $\to$ cascade $\to$ phase-locked $\to$ cascade $\to$ turbulent $\to$ coherent---analogous to the ultradian processing cycle observed in biological coupled oscillator networks.
    (\textbf{D})~Three-dimensional visualisation of R-C-L mode selection on the unit simplex $\Sk + \St + \Se = 1$. The simplex surface is coloured by the dominant processing mode: red (resistive, $\Sk$ minimal), green (capacitive, $\St$ minimal), blue (inductive, $\Se$ minimal). Each point on the simplex represents a unique S-entropy state; the colour indicates which of the three circuit modes the system selects by S-entropy minimisation. The three-fold symmetry reflects the permutation equivalence of the S-entropy coordinates, broken only by the current state of the processing element.}
    \label{fig:processing_cycle}
\end{figure*}

\begin{figure*}[!htbp]
    \centering
    \includegraphics[width=\textwidth]{figures/panel_exp05_performance.png}
    \caption{\textbf{Performance analysis: energy efficiency, throughput, computing stack validation, and self-consistency scaling.}
    (\textbf{A})~Energy per elementary operation on a logarithmic scale for three substrates: silicon CMOS ($\sim 10^{-15}$~J per transition), the Landauer thermodynamic limit $\kB T \ln 2 \approx 2.97 \times 10^{-21}$~J at $T = 310$~K, and the BMD membrane transistor ($\sim 6 \times 10^{-20}$~J per operation). The annotated $10^5 \times$ gap between silicon and the BMD confirms that membrane-mediated computation operates within $20\times$ of the Landauer bound, compared to $10^6 \times$ for current silicon technology.
    (\textbf{B})~Computational throughput (operations per second) on a logarithmic scale for a modern GPU ($\sim 10^{14}$~FLOPS), a $1$~cm$^2$ membrane patch ($\sim 10^{20}$~ops/s), and a $1$~m$^2$ membrane surface ($\sim 10^{24}$~ops/s). Each lipid executes $\sim 10^{11}$ conformational oscillations per second; a $1$~cm$^2$ patch contains $\sim 10^9$ lipids, yielding $\sim 10^{20}$ total partition operations---six orders of magnitude beyond the most powerful binary processors.
    (\textbf{C})~Three-dimensional bar chart of the full computing stack validation across six layers: P-N junction (rectification $47.8\times$), BMD transistor (on/off ratio $42.1$), logic gates (AND/OR/XOR $100\%$ agreement), quantum gates ($88.4\%$ fidelity), gear interconnects ($23{,}500\times$ speedup), and Fibonacci program execution ($91\%$ reliability on a $240$-node scale-free BMD network). All six layers pass their respective validation criteria, confirming bottom-up functional completeness from semiconductor physics to programmable computation.
    (\textbf{D})~Self-consistency deviation (maximum percentage error in closed frequency loops) as a function of molecular complexity (number of atoms) for six benchmark molecules: H$_2$ ($0.00\%$), CO ($0.00\%$), H$_2$O, CO$_2$, CH$_4$ ($0.74\%$), and C$_6$H$_6$ ($1.66\%$). The linear trend line (dashed) shows graceful scaling of error with molecular complexity, consistent with the logarithmic growth predicted by the Compression Theorem. The sub-$2\%$ deviation for the most complex molecule (benzene, $12$ atoms) confirms that the partition framework maintains quantitative accuracy across the molecular hierarchy.}
    \label{fig:performance}
\end{figure*}
